
\documentclass{article}
\usepackage{colortbl}
\usepackage{makecell}
\usepackage{multirow}
\usepackage{supertabular}

\begin{document}

\newcounter{utterance}

\centering \large Interaction Transcript for game `cladder', experiment `full\_v1.5\_default', episode 4979 with Qwen3.6{-}27B.
\vspace{24pt}

{ \footnotesize  \setcounter{utterance}{1}
\setlength{\tabcolsep}{0pt}
\begin{supertabular}{c@{$\;$}|p{.15\linewidth}@{}p{.15\linewidth}p{.15\linewidth}p{.15\linewidth}p{.15\linewidth}p{.15\linewidth}}
    \# & \multicolumn{2}{c}{Player} && \multicolumn{2}{c}{Game Master} \\
    \hline

    \theutterance \stepcounter{utterance}  
    & & & \multicolumn{4}{p{0.6\linewidth}}{
        \cellcolor[rgb]{0.9,0.9,0.9}{
            \makecell[{{p{\linewidth}}}]{
                \texttt{\tiny{[P1$\langle$GM]}}
                \texttt{You are an expert in causal inference. The following question is not a typical commonsense query, but rather a meticulously designed question created by a professor specializing in causal inference, intended to assess the students' mastery of the course content.} \\
\texttt{Imagine a self{-}contained, hypothetical world with only the following conditions, and without any unmentioned factors or causal relationships: Vaccination status has a direct effect on vaccination reaction and getting smallpox. Getting smallpox has a direct effect on black hair. Vaccination reaction has a direct effect on black hair. For unvaccinated individuals, the probability of black hair is 27\%. For vaccinated individuals, the probability of black hair is 37\%. Will vaccination decrease the chance of black hair?} \\
\texttt{Start your answer with 'Yes' or 'No', followed by additional reasoning or evidence to support your explanation.} \\
            }
        }
    }
     \\ \\

    \theutterance \stepcounter{utterance}  
    & \multicolumn{4}{p{0.6\linewidth}}{
        \cellcolor[rgb]{0.561,0.737,0.561}{
            \makecell[{{p{\linewidth}}}]{
                \texttt{\tiny{[P1$\rangle$GM]}}
                \texttt{No} \\
\\ 
\texttt{**Reasoning:**} \\
\\ 
\texttt{To determine whether vaccination *causes* a decrease in the chance of black hair, we must distinguish between the **observed association** (marginal probability) and the **causal effect**. The problem provides marginal probabilities but does not provide the specific conditional probabilities required to calculate the causal effect (such as the probability of black hair given vaccination status and reaction status, or the probability of getting smallpox given vaccination status).} \\
\\ 
\texttt{However, we can analyze the causal structure provided:} \\
\texttt{1.  **Causal Graph:**} \\
\texttt{    *   Vaccination (\$V\$) \$\textbackslash{}rightarrow\$ Vaccination Reaction (\$R\$)} \\
\texttt{    *   Vaccination (\$V\$) \$\textbackslash{}rightarrow\$ Smallpox (\$S\$)} \\
\texttt{    *   Smallpox (\$S\$) \$\textbackslash{}rightarrow\$ Black Hair (\$B\$)} \\
\texttt{    *   Vaccination Reaction (\$R\$) \$\textbackslash{}rightarrow\$ Black Hair (\$B\$)} \\
\\ 
\texttt{    This structure implies that \$V\$ affects \$B\$ through two pathways:} \\
\texttt{    *   Path 1: \$V \textbackslash{}rightarrow S \textbackslash{}rightarrow B\$} \\
\texttt{    *   Path 2: \$V \textbackslash{}rightarrow R \textbackslash{}rightarrow B\$} \\
\\ 
\texttt{2.  **Observed Data:**} \\
\texttt{    *   \$P(B | V=0) = 0.27\$ (Unvaccinated)} \\
\texttt{    *   \$P(B | V=1) = 0.37\$ (Vacc} \\
            }
        }
    }
    & & \\ \\

    \theutterance \stepcounter{utterance}  
    & & & \multicolumn{2}{p{0.3\linewidth}}{
        \cellcolor[rgb]{0.9,0.9,0.9}{
            \makecell[{{p{\linewidth}}}]{
                \texttt{\tiny{[GM$|$GM]}}
                \texttt{No} \\
            }
        }
    }
    & & \\ \\

    \theutterance \stepcounter{utterance}  
    & & & \multicolumn{2}{p{0.3\linewidth}}{
        \cellcolor[rgb]{0.9,0.9,0.9}{
            \makecell[{{p{\linewidth}}}]{
                \texttt{\tiny{[GM$|$GM]}}
                \texttt{no} \\
            }
        }
    }
    & & \\ \\

    \theutterance \stepcounter{utterance}  
    & & & \multicolumn{2}{p{0.3\linewidth}}{
        \cellcolor[rgb]{0.9,0.9,0.9}{
            \makecell[{{p{\linewidth}}}]{
                \texttt{\tiny{[GM$|$GM]}}
                \texttt{game\_result = WIN} \\
            }
        }
    }
    & & \\ \\

\end{supertabular}
}

\end{document}
